\begin{figure}[htbp]
\centering
\begin{tikzpicture}[x=1cm,y=1cm,font=\small,>=Stealth,
  retained/.style={draw=ymAccent,line width=.8pt,rounded corners=2pt,
                   inner xsep=5pt,inner ysep=4pt}]
  \node[anchor=east] at (-1.25,2.05) {Period};
  \node[anchor=east] at (-1.25,1.05) {Spacing $a_N$};
  \node[anchor=east] at (-1.25,-.49) {Spacing $a_{N+1}$};
  \node[anchor=east] at (-1.25,-1.78) {Full $N\to\infty$};
  \node[anchor=east] at (-1.25,-2.35) {Torus kernel};
  % Circles are one periodic coordinate. The number of sites is
  % proportional to the circumference at each of the two resolutions.
  \foreach \cx/\radius/\nsites/\idx in
     {0/.35/8/j,3.05/.525/12/{j+1},6.25/.70/16/{j+2}}{
    \node at (\cx,2.05) {$R_{p(\idx)}$};
    \foreach \cy/\mult in {1.05/1,-.49/3}{
      \begin{scope}[shift={(\cx,\cy)}]
        \draw[ymGrid,line width=.55pt] (0,0) circle[radius=\radius cm];
        \pgfmathtruncatemacro{\total}{\nsites*\mult}
        \pgfmathtruncatemacro{\last}{\total-1}
        \foreach \i in {0,...,\last}{
          \fill[ymInk] ({90+(360/\total)*\i}:\radius) circle[radius=.85pt];}
      \end{scope}}
    \draw[->,ymInk,line width=.65pt] (\cx,-1.40) -- (\cx,-2.03);
  }
  \foreach \cy in {2.05,1.05,-.49,-2.35}
    \node at (8.25,\cy) {$\cdots$};
  \node[retained] (first) at (0,-2.35) {$M^{\rm UV}_{p(j),v}$};
  \node at (3.05,-2.35) {$M^{\rm UV}_{p(j+1),v}$};
  \node[retained] (later) at (6.25,-2.35) {$M^{\rm UV}_{p(j+2),v}$};
  \draw[ymAccent,line width=.8pt] (first.south) -- (0,-3.08);
  \draw[ymAccent,line width=.8pt] (later.south) -- (6.25,-3.08);
  \draw[->,ymAccent,line width=.8pt] (0,-3.08) -- (8.5,-3.08)
    node[right,text=black] {$S_n(v;\cdot)$};
  \node[anchor=north] at (3.10,-3.11) {$k\to\infty$ along $P(k)=p(q(k))$};
\end{tikzpicture}
\caption[Full cutoff limits and a common volume subsequence]{The two
limits for separated correlations. Each column shows one periodic
coordinate at successive lattice spacings (schematic). The full cutoff
sequence at that period gives the torus kernel and the limit for every
separated test supported inside its period cell.
The framed kernels illustrate a subsequence of the prescribed volumes.
One such choice gives convergence in $C^\infty_{\rm loc}(\Conf_n)$ for
all arities and label lists, hence for every separated test function.}
\label{ch13:fig-ordered-limits}
\end{figure}
